\chapter{\IfLanguageName{dutch}{Dataverzameling en voorbereiding}{Data collection and preparation}}%
\label{ch:dataverzameling}

In dit hoofdstuk concentreren we ons op de gebruikte data en hoe deze verzameld en voorbereid zijn voor analyse.
Eerst wordt besproken waar de gegevens vandaan komen, hoe de dataset eruitziet en waarom deze gegevens waardevol zijn voor de bachelorproef.

\section{Dataverzameling}

Voordat de dataset verwerkt kan worden voor de modellen, moet deze eerst verzameld worden.
Daarom werd er vooraf grondig gezocht naar verschillende bronnen met bruikbare ruwe data,
die later tot een gestructureerd geheel konden worden gebracht.
Deze gegevens bestonden uit wedstrijdstatistieken van de Jupiler Pro League,
evenals externe factoren zoals weerdata en stadioninformatie. Nadien werden deze gestructureerd en opgeslagen in een relationele databank.

\subsection{Wedstrijddata}

Eerst werden 7 CSV-bestanden gedownload~\autocite{football_data_belgium}.
Elk bestand bevatte een bepaald seizoen van de Jupiler Pro League, beginnend bij het seizoen 2019-2020 en eindigend vóór de play-offs van het seizoen 2025-2026.
Voorgaande seizoenen werden niet opgenomen omdat er in deze bestanden bepaalde statistieken ontbraken die van belang waren.
Uit deze bron werden diverse wedstrijdstatistieken verzameld,
waaronder de deelnemende teams, de speeldatum en de uiteindelijke winnaar.
Daarnaast werden ook het aantal doelpunten, schoten, overtredingen, corners en kaarten voor zowel de thuis- als uitploeg verzameld.
Voor de feature importance-analyse en de voorspellingen waren deze variabelen uiteraard cruciaal.
Variabelen zoals betting odds werden bewust niet opgenomen omdat deze juist voorspeld moesten worden.
Overige data van de scheidsrechter, teamopstellingen, blessures en transfers waren enkel beschikbaar via betaalde bronnen en werden daarom niet opgenomen.

\subsection{Externe factoren}

Eveneens werden externe factoren toegevoegd aan de dataset.
Hiertoe behoren onder andere weergegevens zoals temperatuur, neerslag, luchtvochtigheid, windsnelheid, bewolking, weercode en luchtdruk.
Deze variabelen werden per uur opgehaald via de API van~\textcite{open_meteo_api}.
Ook deze gegevens waren van belang omdat ze mogelijk een invloed konden hebben op de wedstrijdstatistieken en dus ook op de voorspellingen.

\vspace{5pt}

Aanvullend werden de coördinaten van de verschillende stadions opgezocht. Deze waren vereist als parameters voor het ophalen van de weerdata.
Naast de ligging van het stadion werd ook de stadioncapaciteit opgenomen als feature. 
Hoewel deze variabele geen exacte maat vormt voor het effectieve aantal aanwezige toeschouwers, kan ze wel dienen als een benadering van de schaal van het thuispubliek en de infrastructuur van de club. 
Zo kon onderzocht worden of stadiongrootte en het aantal supporters, met merendeels steun voor de thuisploeg, samenhangt met bepaalde wedstrijdstatistieken.

\subsection{Dataopslag}

Na het verzamelen van de verschillende gegevens werden deze gestructureerd en opgeslagen in een relationele databank om consistentie en toegankelijkheid te garanderen.
Deze databank werd beheerd met behulp van Microsoft SQL Server in SQL Server Management Studio (SSMS).
Figuur \ref{fig:databank-schema} toont de structuur van de databank met de verschillende tabellen en relaties.

\begin{figure}[h!]
\centering
\includegraphics[width=0.8\textwidth]{databank_schema.png}
\caption{Schematische weergave van de relationele databank}
\label{fig:databank-schema}
\end{figure}

\vspace{5pt}

Hierin werd alle data van de wedstrijdstatistieken in de tabel \textit{WedstrijdStatistiek} opgeslagen.
Deze tabel verwees naar de tabel \textit{Wedstrijd}, waarin de details van elke wedstrijd werden bijgehouden, zoals datum, tijd en teams.
Daarnaast was er nog een tabel \textit{Team}, die de locatie en capaciteit van elk stadion bevatte.
Overigens waren er nog twee tabellen voor de weerdata. De tabel \textit{WedstrijdWeer} bevatte de weersomstandigheden op het moment van de aftrap. 
Omdat de API weergegevens per uur aanleverde, werd het uur van de wedstrijd als referentiepunt genomen. 
Daarnaast werden ook gemiddelden berekend binnen een tijdsvenster van één uur rond de aftrap om de weersomstandigheden tijdens de wedstrijd beter te benaderen.
De tabel \textit{WedstrijdWeerVooraf} bevatte weergegevens van de 48 uur voorafgaand aan de wedstrijd. 
Deze bevatten het gemiddelde van de temperatuur, de som van de hoeveelheid neerslag en de maximale windsnelheid.